\section{Georges Saab}
Senior Vice President do Java Platform Group na Oracle e chair do OpenJDK
Governing Board; na JavaSoft foi fundador do grupo do Swing e do
Java Web Start, e depois dirigiu a JRockit na BEA (fonte:
\citeobra{bio:saab-inside}).

\begin{itemize}
\obra{video:saab-rapid-innovation}{Por que o Java passou a lançar uma
  versão a cada seis meses (GOTO 2018).}
\obra*{video:saab-java-2024}{Keynote sobre o estado da plataforma em 2024.}
\obra*{video:saab-podcast-43}{Podcast: previsibilidade, que problemas o
  Java deve atacar e o foco recente em estudantes.}
\end{itemize}
Material escrito por ele é pouco; as palestras acima são o que há de mais
útil.

\section{Sharat Chander}
Hoje aparece como \enquote{Oracle Alum} no inside.java (fonte:
\citeobra{bio:chander-inside}); o Foojay o descreve como ex-Senior Director
de Java Product Management e Developer Engagement na Oracle.

\begin{itemize}
\obra{artigo:arrival-java21}{O anúncio oficial do JDK~21~LTS, com a lista
  de JEPs.}
\obra{artigo:java30-faceoff}{Retrospectiva dos 30 anos, com os recursos
  mais votados pela comunidade.}
\obra*{video:chander-25th}{Entrevista sobre a comunidade Java nos 25 anos.}
\end{itemize}

\section{Heather VanCura}
VP de Community Engagement na Oracle, à frente do programa Java Community
Process (JCP) (fonte: \citeobra{bio:vancura-qcon}).

\begin{itemize}
\obra{video:vancura-productivity}{Como adotar o Java atual e contribuir com
  a evolução da plataforma: JCP e grupos de usuários (JUGs).}
\obra*{artigo:vancura-qa}{Entrevista sobre o primeiro ano dela como chair do
  JCP e o papel das especificações.}
\obra{doc:jcp}{Onde ficam públicos os JSRs do Java~SE e os grupos de
  trabalho; qualquer pessoa pode acompanhar.}
\end{itemize}
